\textbf{Время:} 01:45

\textbf{Тезис:} корреляция интенсивностей может нести информацию об угловом размере источника.

\textbf{Визуал:} звезда → два приёмника → сравнение сигналов при разных расстояниях между приёмниками.

\textbf{Текст:} Звезда далека, и её угловой размер мал. Хэнбери Браун и Твисс предложили получать сведения о размере через связь между двумя зарегистрированными сигналами. В 1956 году после лабораторной демонстрации они применили интенсивностную интерферометрию к Сириусу.

Здесь сравниваются флуктуации фототоков двух раздельных светоприёмников. Каждый приёмник получает свет от всей звезды; стрелки на схеме показывают каналы регистрации, а не отдельные участки диска. Меняя расстояние между приёмниками, исследуют, как меняется корреляция интенсивностей. Её пространственный масштаб позволяет оценить угловой размер звезды. Это не обычное изображение диска и не мерцание звезды, которое видно глазом.

Не нужно сейчас осваивать весь оптический расчёт. Запомним конструкцию вопроса: один приёмник сообщает, сколько света пришло; два дают ещё и информацию о совместной регистрации. В этой дополнительной информации обнаруживается размер источника.

Пауза: показать, что меняется в опыте — расстояние между приёмниками, а не сама звезда.

\textbf{Переход:} Как общий принцип парных корреляций проявился у частиц в лаборатории?
